\appendix
\setcounter{theorem}{0}
\renewcommand{\thetheorem}{A.\arabic{theorem}}
\renewcommand{\theequation}{A.\arabic{equation}}
\renewcommand{\thesubsection}{A.\arabic{subsection}}

\section*{Appendix}

\subsection{ODEs in Banach spaces} We now introduce the two different notions of ODEs in Banach spaces used in the present article. Let us consider the framework of \cref{sec:abstract-construction} and let $I\subset \R$ be a nonempty interval and take $s_0\in I$.

For any $s\in \R$, we~can easily extend $e^{sA}$ to $C^{0}([0,T],X^{\sigma})$ by the formula
\begin{align*}
\left[e^{sA}V\right](t)=e^{sA}V(t),\quad \text{for $V\in C^{0}([0,T],X^{\sigma})$.}
\end{align*}
Given $H\in L^{1}\big(I,C^{0}([0,T],X^{\sigma})\big)$, we~say that $\xi\in C^{0}\big(I,C^{0}([0,T],X^{\sigma})\big)$ satisfies
\begin{align}\label{appendix:ode:ode-banach-1}
\left\{\begin{aligned}
\dfrac{d}{ds}\xi(s)&=A \xi(s)+H(s), && s\in I, \\
\xi(s_0)&=\xi_0, &&
\end{aligned}\right.
\end{align}
with $\xi_0\in C^{0}([0,T],X^{\sigma})$, if~it satisfies
\begin{align}\label{appendix:ode:ode-banach-1-duhamel}
\xi(s)&=e^{(s-s_0)A}\xi_0+\int_{s_0}^{s}e^{(s-w)A}H(w)dw,\ \forall s\in I,
\end{align}
with equality in $C^{0}([0,T],X^{\sigma})$.
\begin{lemma}
\label{lmDuhamelCauchy}
If $H\in C^0\left(I,C^{0}([0,T],\P_{n}X^{\sigma})\right)$ and $\xi\in C^{0}\left(I,C^{0}([0,T],\P_{n}X^{\sigma})\right)$ for some $n\in\N$ and satisfies
$ \frac{d}{ds}\xi(s)=A \xi(s)+H(s)$ in the previous sense \eqref{appendix:ode:ode-banach-1-duhamel}, then $\xi\in C^{1}\left(I,C^{0}([0,T],\P_{n}X^{\sigma})\right)$ and is a classical solution.
\end{lemma}
\begin{proof}
Since $A$ is a bounded operator on $C^{0}([0,T],\P_{n}X^{\sigma})$, $\frac{d}{dt}e^{tA}\xi_0=A e^{tA}\xi_0$ in the classical sense of $C^1$ functions with value in $C^{0}([0,T],\P_{n}X^{\sigma})$. In~particular, Duhamel's formula \eqref{appendix:ode:ode-banach-1-duhamel} leads to the claimed result by usual arguments of semigroup theory.
\end{proof}

\begin{lemma}\label{lmtranslat}
For $T_1<T_2$, let us consider $T\in (0, T_2-T_1)$, $\eta\in (0, T_2-T-T_1)$ and $I:=[T_1-\eta, T_2-T-\eta]$. Let $G\in C^{0}([T_{1},T_{2}],X^{\sigma})$ and assume that $V\in C^{0}([T_{1},T_{2}],X^{\sigma})$ is a mild solution of
\begin{align*}
\left\{\begin{aligned}
\dfrac{d}{dt}V(t)&=A V(t)+G(t) &&\text{ for }\ t\in [T_1, T_2],\\
V(T_1)&=V_0. &
\end{aligned}\right.
\end{align*}
If we define $\xi,H\in C^{0}\bigl(I, C^{0}([0,T],X^{\sigma})\bigr)$ by $\xi(s)=V^s$ and $H(s)=G^s$ with
\begin{align*}
\xi(s)(t)=V^{s}(t)=V(t+s+\eta)\ \text{ and }\ H(s)(t)=G^{s}(t)=G(t+s+\eta),
\end{align*}
for all $s\in I, t\in [0,T]$, then for any $s_0\in I$, $\xi$ is solution, in the sense of \eqref{appendix:ode:ode-banach-1-duhamel}, of
\begin{align*}
\left\{\begin{aligned}
\dfrac{d}{ds}\xi(s)&=A \xi(s)+H(s), && s\in I, \\
\xi(s_0)&=\xi_0, &
\end{aligned}\right.
\end{align*}
with $\xi_0=V^{s_0}=V(\cdot+s_0+\eta)$.
\end{lemma}
\begin{proof}
By Duhamel's formula, for all $t\in [T_1, T_2]$ we have
\begin{align*}
V(t)&=e^{(t-T_1)A}V_{0}+\int_{T_1}^{t}e^{(t-\tau)A}G(\tau)d\tau,
\end{align*}
with $V(T_1)=V_0$. Pick $s_0\in I$. First observe that, for any $t\in [0, T]$,
\begin{align*}
V(t+s_0+\eta)=e^{(t+s_0+\eta-T_1)A}V_{0}+\int_{T_1}^{t+s_0+\eta}e^{(t+s_0+\eta-\tau)A}G(\tau)d\tau.
\end{align*}
Then for $s\in I$ and $t\in [0, T]$,
\begin{align*}
V^{s}(t)&=V(t+s+\eta)=e^{(t+s+\eta-T_1)A}V_{0}+\int_{T_1}^{t+s+\eta}e^{(t+s+\eta-\tau)A}G(\tau)d\tau\\
&=e^{(s-s_0)A}V(t+s_0+\eta)+\int_{s_0}^{s}e^{(s-w)A}G(t+w+\eta)dw\\
&=e^{(s-s_0)A}V^{s_0}(t)+\int_{s_0}^{s}e^{(s-w)A}G^{w}(t)dw.
\end{align*}
So, since this is true for any $t\in [0,T]$, it~gives
\begin{align*}
V^{s}=e^{(s-s_0)A}V^{s_0}+\int_{s_0}^{s}e^{(s-w)A}G^{w}dw,\quad \forall s\in I.
\end{align*}
By hypothesis, this equality holds in $C^0([0, T], X^\sigma)$, and is exactly \eqref{appendix:ode:ode-banach-1-duhamel}, as we wanted to prove.
\end{proof}
